\section{Results}
\label{sec:results}

We evaluate CT-KAT on seven harness rows across ML-KEM, ML-DSA, and synthetic
controls, under up to six gcc/clang build cells. The committed corpus CSVs record
build-cell provenance. Table~\ref{tab:coverage} shows each layer catches a case
the others miss; Table~\ref{tab:corpus} grounds verdict classes; Table~\ref{tab:dudect}
is the timing appendix. Tables~\ref{tab:ablation} and~\ref{tab:mldsa} add the
two supporting checks: what each omitted layer misses, and why ML-DSA remains
\texttt{needs-analysis}.

\subsection{Per-feature single-coverage}
\label{sec:results-coverage}

Table~\ref{tab:coverage} names, per layer, one case that \emph{only} it flags.
Two are the paper's spine.

\paragraph{(i) asm-scan: structural CT alone is insufficient.}
The KyberSlash harness passes Memcheck clean: no secret-tainted branch or
address. Yet asm-scan flags a
secret-derived integer division---the KyberSlash class~\cite{kyberslash}---whose
latency is data-dependent. Triage marks it secret-derived, giving
the risk verdict in Table~\ref{tab:corpus}; a Memcheck-only deployment would
ship it as clean.

\paragraph{(ii) Default-deny triage: a structural FAIL the tool refused to
call a leak.}
On real ML-DSA-65 (\texttt{mldsa65} / \texttt{sign}) structural checking fails.
Three --- \texttt{poly\_chknorm}, \texttt{make\_hint}, and
\texttt{poly\_challenge} --- are registered rejection-sampling behavior required
by ML-DSA~\cite{fips204,dilithium}. Two are not:
\texttt{crypto\_sign\_signature\_ctx} (optimized-build inlining) and
\texttt{pack\_sig} (a public-signature packer). Default-deny therefore holds the
row at \texttt{needs-analysis}: neither auto-accepting nor over-claiming. The
\texttt{-O0 -fno-inline} cells show only registered functions; optimized cells
surface the two unregistered names (Table~\ref{tab:mldsa}).

\paragraph{The remaining three layers.}
The synthetic controls justify the remaining layers. \texttt{toy\_lookup/leaky}
is a secret-indexed lookup Memcheck flags in all six cells. \texttt{ct\_matrix\_flip/leaky}
fails at \texttt{-O0} but passes at \texttt{-O2}/\texttt{-Os}, exposing
build-sensitive CT. \texttt{toy\_dudect} isolates a timing leak with no structural
branch. No four layers subsume the fifth; Table~\ref{tab:ablation} states the
corresponding miss if each layer is removed.

\subsection{Verdict-class corpus}
\label{sec:results-corpus}

Table~\ref{tab:corpus} grounds every verdict class in a concrete target. Stock
\texttt{mlkem768/kem\_dec}~\cite{pqclean,fips203,kyber} is \texttt{robust}; its
Keccak division candidates triage public. The KyberSlash variant is
\texttt{varlat-secret-risk}; ML-DSA contributes \texttt{needs-analysis}; and
synthetic controls cover \texttt{ct-leak} and \texttt{build-sensitive-ct}. The
tool never auto-declares \texttt{ct-leak}: an explicit human verdict is required.
Together the tables make the framework layer-justified in this corpus and
validated on concrete runs, without claiming formal completeness or leak absence.

\subsection{Statistical timing appendix and its honesty caveat}
\label{sec:results-dudect}

dudect is keyed per target, not per build cell, so it lives beside
Table~\ref{tab:corpus}. The \texttt{toy\_dudect}
rows are unambiguous: \texttt{leaky} ($n_0=\dudectLeakyNzero$,
$n_1=\dudectLeakyNone$, mean cycles $\dudectLeakyMeanZero$ versus
$\dudectLeakyMeanOne$) gives $|t|=\dudectLeakyT$ and a FAIL; \texttt{safe}
gives $|t|=\dudectSafeT$ and a PASS.

We report the caveat because the tool does: the \texttt{leaky} run dropped
$\dropTotalPct\%$ zero-cycle samples, asymmetrically
($\dropClassZeroPct\%$ of class 0 versus $\dropClassOnePct\%$ of class 1), from
QEMU/Docker TSC skew. The large separation ($|t|=\dudectLeakyTrunc$) makes the
FAIL unambiguous, but publication-grade timing should be confirmed natively. The
same warning appears on stock \texttt{mlkem768} ($t=\mlkemDudectT$), read as QEMU
noise rather than a real leak.
